Let $n\in\mathbb N$ and let $A\subseteq\{1,2,\dots,n\}$. Define a simple graph $G(A)$ with vertex set $A$ and an (undirected) edge $\{x,y\}$ for distinct $x,y\in A$ iff $\gcd(x,y)=1$.
Define
\[T(n):=\Big\lfloor\frac n2\Big\rfloor+\Big\lfloor\frac n3\Big\rfloor-\Big\lfloor\frac n6\Big\rfloor.
\]

Question 1: Is it true that if $|A|>T(n)$ then $G(A)$ contains, for every odd integer $L$ with $3\le L\le \lfloor n/3+1\rfloor$, a (simple) cycle of length $L$ as a subgraph?

Question 2: For each fixed $\ell\ge 1$, is it true that for all sufficiently large $n$, every $A\subseteq\{1,\dots,n\}$ with $|A|>T(n)$ has $G(A)$ containing a complete tripartite subgraph with parts of sizes $(1,\ell,\ell)$ (denoted $K_{1,\ell,\ell}$), i.e. $2\ell+1$ vertices with all edges between different parts present?
